% ─────────────────────────────────────────────────────────────────────────────
% LabBank — one-shot ALM system report (10-12 pages)
% Compile: pdflatex main.tex  (run twice for the table of contents)
%
% STATUS: first-iteration PLAN, not a finished camera-ready document.
%
% Done: pipeline.png (post-Dagster-bracket, exported from the draw.io GUI) is
% in images/; the balance-sheet-tab annotated mockup is also in images/, as a
% plain PNG (labbank_balance_sheet_annotated.png) — it used to be embedded as
% an .svg via the svg package/Inkscape, but that pipeline mis-rendered the
% text overlay (garbled/duplicated glyphs) on Overleaf's compiler, so it was
% rasterised once to PNG and embedded with \imgfull like every other figure.
%
% Remaining TODOs before this is publish-ready:
%   1. Replace the two \sqlplaceholder boxes in Section 7 ("See It Yourself")
%      with real SSMS/Azure Data Studio screenshots.
%   2. Replace the two \figplaceholder "before/after" boxes in Section 3
%      ("Change Something, See What Happens") with real before/after
%      screenshots of the Metrics tab.
%   3. Fill in the LinkedIn / email placeholders — title page + "Contact"
%      section at the end (search REPLACE_WITH_).
%   4. Update \reportdate below if you want something other than today's date.
% ─────────────────────────────────────────────────────────────────────────────
\documentclass[11pt,a4paper]{article}

\usepackage[a4paper,top=2.2cm,bottom=2.4cm,left=2.3cm,right=2.3cm]{geometry}
\usepackage[utf8]{inputenc}
\usepackage[T1]{fontenc}
\usepackage{graphicx}
\usepackage{xcolor}
\usepackage{booktabs}
\usepackage{longtable}
\usepackage{array}
\usepackage{enumitem}
\usepackage{microtype}
\usepackage{titlesec}
\usepackage{fancyhdr}
\usepackage[most]{tcolorbox}
\usepackage{amsmath}
\usepackage{hyperref}

\graphicspath{{images/}}

% ── Palette (matches visual_rep/labbank_presentation.tex) ──────────────────────
\definecolor{labblue}{RGB}{21,  101, 192}
\definecolor{labred} {RGB}{198,  40,  40}
\definecolor{labgrn} {RGB}{ 46, 125,  50}
\definecolor{labpur} {RGB}{106,  27, 154}
\definecolor{labdark}{RGB}{ 26,  35, 126}
\definecolor{labgrey}{RGB}{ 96, 125, 139}

\newcommand{\reportdate}{\today}

% ── Section styling ─────────────────────────────────────────────────────────────
\titleformat{\section}
  {\normalfont\Large\bfseries\color{labdark}}
  {\thesection}{0.6em}{}[{\color{labblue!40}\titlerule[1pt]}]
\titleformat{\subsection}
  {\normalfont\large\bfseries\color{labblue}}
  {\thesubsection}{0.6em}{}

\pagestyle{fancy}
\fancyhf{}
\renewcommand{\headrulewidth}{0.4pt}
\renewcommand{\footrulewidth}{0pt}
\fancyhead[L]{\small\color{labgrey}LabBank --- Laboratory Banking ALM System}
\fancyhead[R]{\small\color{labgrey}\thepage}
\fancypagestyle{plain}{\fancyhf{}\renewcommand{\headrulewidth}{0pt}}

\hypersetup{
  colorlinks=true,
  linkcolor=labblue,
  urlcolor=labblue,
  citecolor=labblue,
  pdftitle={LabBank -- A Laboratory Banking ALM System},
  pdfauthor={LabBank Project}
}

% ── Convenience macros ───────────────────────────────────────────────────────────
\newcommand{\imgfull}[2][0.92\textwidth]{%
  \begin{center}\includegraphics[width=#1]{#2}\end{center}}
\newcommand{\pos}[1]{\textcolor{labgrn}{\textbf{#1}}}
\newcommand{\neg}[1]{\textcolor{labred}{\textbf{#1}}}
\newcommand{\hi} [1]{\textcolor{labblue}{\textbf{#1}}}

\newtcolorbox{pitchbox}[1][Why this matters]{
  enhanced, breakable, colback=labblue!5, colframe=labblue!55,
  boxrule=0.8pt, arc=2mm, left=8pt, right=8pt, top=6pt, bottom=6pt,
  title={#1}, colbacktitle=labdark, coltitle=white, fonttitle=\bfseries\small,
  attach boxed title to top left={xshift=8pt, yshift=-2mm},
  boxed title style={boxrule=0pt, arc=1mm}
}

\newtcolorbox{sqlplaceholder}[2][7.2cm]{
  enhanced, breakable, colback=white, colframe=labgrey,
  boxrule=1pt, sharp corners=south, arc=1mm,
  title={#2}, colbacktitle=labdark, coltitle=white, fonttitle=\bfseries\small,
  height=#1, valign=center, halign=center
}

% Generic "asset not ready yet" box — dashed border so it reads as a plan
% placeholder, not a broken figure. #1 = height, #2 = what belongs here.
% Currently unused (both pipeline.png spots now call \imgfull directly) —
% kept in case images/pipeline.png isn't ready yet when you compile; swap
% back to \figplaceholder[7cm]{...} at either \imgfull[...]{pipeline.png} call.
\newtcolorbox{figplaceholderbox}[1][6cm]{
  enhanced, breakable, colback=labgrey!4, colframe=labgrey,
  boxrule=0pt, arc=1mm, borderline={1pt}{-1pt}{labgrey, dashed},
  height=#1, valign=center, halign=center
}
\newcommand{\figplaceholder}[2][6cm]{%
  \begin{figplaceholderbox}[#1]
    \normalfont\color{labgrey}
    {\bfseries PLACEHOLDER --- not yet exported}\\[4pt]
    \itshape #2
  \end{figplaceholderbox}
}

\newcommand{\schemarow}[4]{\texttt{#1} & \texttt{#2} & #3 & #4 \\}

% ─────────────────────────────────────────────────────────────────────────────
\begin{document}
% ─────────────────────────────────────────────────────────────────────────────

% ═══ TITLE PAGE ═══
\begin{titlepage}
\begin{center}
\vspace*{0.6cm}
{\color{labdark}\fontsize{40}{44}\selectfont\bfseries LabBank}\\[8pt]
{\color{labgrey}\Large A Laboratory Banking ALM System}\\[22pt]
{\large\color{black} A digital twin of a bank's balance sheet and transaction ledger ---\\[2pt]
 built to be explored, stress-tested, and rebuilt with your own numbers}\\[22pt]

\imgfull[0.94\textwidth]{pipeline.png}

\vspace{10pt}
{\small\itshape\color{labgrey}The pipeline that builds it (above) and the app you explore it with (below).}

\vspace{8pt}
\imgfull[0.56\textwidth]{labbank_balance_sheet_annotated.png}

\vspace{16pt}
{\small\color{labgrey}
Python \;\textbullet\; SQL Server \;\textbullet\; Dagster \;\textbullet\; Streamlit \;\textbullet\; EBA/RTS/2022/10}\\[6pt]
{\small\color{labgrey}\reportdate}
\end{center}
\vfill
\begin{center}
\small \href{https://github.com/dzimius/LabBank}{github.com/dzimius/LabBank}\quad|\quad
\href{mailto:REPLACE_WITH_YOUR_EMAIL}{REPLACE\_WITH\_YOUR\_EMAIL}\quad|\quad
\href{https://www.linkedin.com/in/REPLACE_WITH_YOUR_HANDLE}{linkedin.com/in/REPLACE\_WITH\_YOUR\_HANDLE}
\end{center}
\end{titlepage}

\tableofcontents
\thispagestyle{plain}
\newpage

% =============================================================================
\section{The Problem, and the Idea}
% =============================================================================

LabBank is an educational, hands-on ALM system: a \textbf{digital twin} of a
bank's balance sheet and transaction ledger, built so you \textbf{learn
interest rate risk and liquidity risk by testing them}, not by reading about
them. Build your own balance sheet, edit a product's rate or a deposit's decay
assumption, and watch the change ripple all the way through to NII, EVE, LCR,
and NSFR --- the same regulatory metrics a real bank reports. Nothing here is a
static example: every product, rate, and cash-flow row is real data you can
inspect, edit, and regenerate.

\textbf{Why that's hard to get right.} A bank's balance sheet produces
thousands of cash flows every day, and rates reset, balances prepay, deposits
decay, and shock scenarios apply simultaneously across every position in every
currency. Most systems hide that complexity behind aggregated dashboards, so
the moment a number looks wrong, nobody can point at the row that produced it,
and nobody can safely ask ``what if'' without a change-request ticket. LabBank
makes the opposite choice on both counts: every metric is something you can
test yourself, and every number is something you can trace yourself.

\begin{pitchbox}[Core principle]
Every number in every report traces back to a \textbf{single row} in a
\textbf{single cash-flow table}. No aggregation happens before it is recorded.
That means NII, EVE, LCR, and NSFR are not four separate models --- they are four
different aggregations of the \textit{same} underlying cash-flow schedule.
\end{pitchbox}

\subsection{What this project demonstrates}
\begin{itemize}[itemsep=2pt]
  \item A \textit{complete} ALM pipeline built from scratch: raw transaction
        generation $\rightarrow$ market/behavioural enrichment $\rightarrow$
        cash flows $\rightarrow$ regulatory metrics $\rightarrow$ reporting.
  \item An interactive sandbox (the LabBank Streamlit app,
        Section~\ref{sec:labbankapp}) where
        editing the balance sheet mix, the IRS book, or an NMD decay curve
        immediately recomputes NII/EVE/SOT/LCR/NSFR --- so you can see, not
        just read, how a given ALM lever moves the numbers.
  \item How NII, EVE, LCR, and NSFR are all derived from the same cash-flow
        rows, just aggregated differently --- so there is no reconciliation gap
        between risk views.
  \item How an interest-rate-swap overlay shifts the repricing gap without
        touching the underlying balance sheet.
  \item How Dagster orchestration enforces dependency order across the pipeline
        and gives every run a queryable audit trail.
  \item A SQL Server data model designed to be inspected directly ---
        \texttt{SELECT * FROM cf.products} is a legitimate way to understand
        this system, not a debugging escape hatch.
\end{itemize}

\subsection{Who this is for}
Two audiences, one codebase. \textbf{Practitioners} who sit between ALM, risk
analytics, finance, data engineering, and reporting, and want to see how those
pieces actually connect end to end. \textbf{Learners} --- students, career
switchers, anyone studying IRRBB/liquidity methodology --- who want a real
balance sheet to experiment on instead of a static textbook example: change an
assumption, break something, see exactly which number moved and why. Either
way, it is not a toy notebook and not a closed black-box model: every formula
has a direct SQL or Python equivalent you can query yourself.

% =============================================================================
\section{Architecture at a Glance}
\label{sec:architecture}
% =============================================================================

Four layers, one integration bus:
\[
\text{Excel inputs} \;\rightarrow\; \text{Python modules (Dagster-orchestrated)}
\;\rightarrow\; \text{SQL Server (9 schemas)} \;\rightarrow\; \text{Reports}
\]

\imgfull[\textwidth]{pipeline.png}

SQL Server is the single integration bus. Every stage reads from tables the
previous stage wrote, instead of passing hidden in-memory state --- which means
every stage can be re-run independently, and every intermediate result can be
audited by querying the database directly. Dagster sits on top of all four
processing stages (balance-sheet generation, SQL storage, IRRBB/liquidity
calculation, and reporting) and enforces the order they run in.

\subsection{Pipeline, step by step}
\begin{enumerate}[itemsep=2pt]
  \item \textbf{Balance sheet generation} --- synthetic transactions drawn from
        \texttt{bank\_data.xlsx} using truncated-normal balance distributions
        and real historical WIBOR rates. Written to \texttt{dbo.transactions}
        and five \texttt{schemat.*} tables.
  \item \textbf{Market data enrichment} --- yield curves bootstrapped into
        \texttt{mkt.curves}; historical fixings into \texttt{mkt.fixings};
        behavioural parameters (prepayment, NMD decay) into \texttt{bs.*}.
  \item \textbf{Cash flow generation} --- for every transaction, a full schedule
        of contractual amortisation plus behavioural adjustments is written to
        \texttt{cf.products}. Repricing gaps land in \texttt{irrbb.ir\_gap\_*}.
  \item \textbf{IR derivatives (optional)} --- swap legs from
        \texttt{irs\_input.xlsx}, cash flows written to \texttt{cf.ir\_swap\_*},
        swap gap merged into \texttt{irrbb.ir\_gap\_beh}.
  \item \textbf{IRRBB calculation} --- six EBA shock scenarios applied to every
        curve; NII over a 1-year horizon, EVE as the present value of run-off
        cash flows.
  \item \textbf{Liquidity calculation} --- LCR and NSFR, written to
        \texttt{results.lcr\_nsfr}.
  \item \textbf{Reporting} --- Jupyter notebooks read live from SQL and export
        to HTML/PDF; Excel outputs for IRRBB and LCR/NSFR; a Streamlit app for
        interactive stress-testing.
\end{enumerate}

% =============================================================================
\section{The LabBank App --- Your Digital Twin, Interactively}
\label{sec:labbankapp}
% =============================================================================

Everything in Section~\ref{sec:architecture} exists so that this app has
something real to sit on top of. The Streamlit app (\texttt{sandbox/app.py}) reads a pre-built demo
balance sheet --- no SQL Server required --- and lets you edit it directly:
change the product mix, the IRS book, or a non-maturity-deposit decay
assumption, and every downstream metric recomputes live.

\begin{center}
\ttfamily\small
python -m venv venv \quad.\textbackslash venv\textbackslash Scripts\textbackslash Activate.ps1\quad
pip install -r requirements.txt\quad streamlit run sandbox\textbackslash app.py
\end{center}

\subsection{The Balance Sheet tab}
The tab you land on. Edit the asset and liability mix directly --- every other
tab reacts to whatever you set here.

\imgfull[0.92\textwidth]{labbank_balance_sheet_annotated.png}

\subsection{The other five tabs}
\begin{itemize}[itemsep=2pt]
  \item \textbf{IRS Book} --- add/remove swap trades, change notional or fixed
        rate, toggle \texttt{pay\_fixed} per swap; a leg summary compares
        baseline vs.\ your edits.
  \item \textbf{NMD Stress} --- pick a non-maturity deposit product and edit its
        decay profile directly; a chart compares your stressed profile against
        the baseline.
  \item \textbf{Metrics} --- the results dashboard: NII, EVE, LCR, NSFR baseline
        vs.\ modified, the EBA Supervisory Outlier Test with pass/fail called
        out explicitly, per-scenario $\Delta$NII/$\Delta$EVE charts, and an IRS
        contribution breakdown.
  \item \textbf{Gap Analysis} --- repricing-gap panels (with and without the IRS
        overlay) and the 12-month behavioural liquidity gap.
  \item \textbf{Market Curves} --- base and shocked curves per currency, plus a
        panel to try rate environments that aren't one of the 7 EBA scenarios.
\end{itemize}

\subsection{Change something, see what happens}
This is the point of the whole exercise. Every edit above is additive and
immediate --- there is no ``recalculate'' button to press and no batch job to
wait on.

\begin{pitchbox}[A worked example]
Open \textbf{Balance Sheet}, shift 10 percentage points of asset mix from
fixed-rate mortgages into floating-rate mortgages, keeping both sides at
100\%. Switch to \textbf{Metrics}: the repricing gap becomes more
asset-sensitive, so a parallel-up shock now helps NII \textit{more} than
before, and a parallel-down shock hurts it more --- while $\Delta$EVE moves the
other way, since floating-rate assets carry far less duration risk than the
fixed-rate mortgages they replaced. Same balance sheet size, same products,
one slider moved --- and the NII-vs-EVE trade-off from
Section~\ref{sec:irrbb} is now something you watched happen, not just a
formula you read.
\end{pitchbox}

\begin{minipage}[t]{0.48\textwidth}
\begin{figplaceholderbox}[6cm]
  \normalfont\color{labgrey}
  {\bfseries Metrics tab --- baseline}\\[4pt]
  \itshape [ screenshot before the edit ]
\end{figplaceholderbox}
\end{minipage}\hfill
\begin{minipage}[t]{0.48\textwidth}
\begin{figplaceholderbox}[6cm]
  \normalfont\color{labgrey}
  {\bfseries Metrics tab --- after shifting 10pp fixed $\rightarrow$ floating}\\[4pt]
  \itshape [ screenshot after the edit ]
\end{figplaceholderbox}
\end{minipage}

\vspace{6pt}
\small\color{labgrey}\textit{Any lever works for this comparison --- NMD decay
speed, swap notional, the balance sheet mix itself. The point is that changing
one input and re-reading the Metrics tab \textit{is} the analysis; there's no
separate model to reconcile it against.}
\normalsize\color{black}

% =============================================================================
\section{Orchestration --- Dagster}
% =============================================================================

Business logic lives in plain Python workflow scripts; Dagster wraps each one
as a software-defined asset. That buys dependency enforcement, retries, run
history, and a live web UI --- without rewriting a single calculation.

\begin{center}
\small
\begin{tabular}{ll}
\toprule
\textbf{Job} & \textbf{Purpose} \\
\midrule
\texttt{balance\_sheet\_job}  & Generate the balance sheet, load enrichment data, process IRS positions. \\
\texttt{full\_run\_job}       & Complete pipeline: BS $\rightarrow$ CFs $\rightarrow$ IRS $\rightarrow$ IRRBB + LCR/NSFR. \\
\texttt{irs\_update\_job}     & Reload swaps, recompute IRRBB only --- balance sheet unchanged. \\
\texttt{irrbb\_recalc\_job}   & Recompute NII / EVE / SOT from existing cash flows after a curve update. \\
\texttt{liq\_only\_job}       & Standalone LCR / NSFR refresh. \\
\texttt{optimize\_prep\_job}  & Rebuild fast-approximation tensors and run the accuracy check. \\
\texttt{labbank\_data\_job}   & Full pipeline + optimize\_prep in one step --- the one to run after editing your own input data. \\
\bottomrule
\end{tabular}
\end{center}

\begin{pitchbox}[Key property]
Steps 5, 6, and 7 of the pipeline can each be re-run independently without
touching the balance sheet or cash flows --- enforced by the Dagster job
definitions, not by convention.
\end{pitchbox}

% =============================================================================
\section{Data Model --- SQL Schema Overview}
% =============================================================================

Nine schemas, each owned by one stage of the pipeline. Downstream stages read
from upstream tables rather than sharing in-memory objects, which is what makes
every stage independently re-runnable and auditable with a plain SQL client.

\imgfull[\textwidth]{sql_schema.png}

If you are exploring the database yourself: \texttt{dbo.transactions} and
\texttt{schemat.*} answer ``what does the bank actually hold'', \texttt{cf.products}
answers ``what cash flows come out of that'', and \texttt{irrbb.irrbb\_report} /
\texttt{results.lcr\_nsfr} give you the final regulatory-style numbers.

\subsection{The core entity relationships}
The diagram above shows how data \textit{flows} between schemas; the one below shows
how the master ledger is \textit{structured} --- the same view you would get by
right-clicking a database in SQL Server Management Studio and choosing ``Database
Diagram''. \texttt{dbo.transactions} is the hub: every position the bank holds is
exactly one row there. Each \texttt{schemat.*} table adds product-specific detail for
that same \texttt{transaction\_id} --- a 1:1 relationship, since a given position is a
loan, a deposit, a security, or an equity item, never more than one at once.

\imgfull[0.9\textwidth]{transaction_erd.png}

% =============================================================================
\section{Schedule of SQL Tables}
\label{sec:sqltables}
% =============================================================================

Nine schemas, roughly thirty tables --- every one of them traceable to the
pipeline stage that writes it (Section~\ref{sec:architecture}) and the ERD
above. The full table-by-table inventory (schema, table, key columns,
purpose) is reference material rather than something a document this length
needs to print in full --- it lives in the ``Full table reference'' section
of the methodology document (\texttt{docs/02\_methodology.md}), referenced in
Section~\ref{sec:tryityourself}.

% =============================================================================
\section{See It Yourself --- Live SQL Server Queries}
\label{sec:seeitforyourself}
% =============================================================================

Every claim in this document can be checked with a \texttt{SELECT} statement
against the running database. Two examples below --- what a single contract
looks like end to end, and what the final regulatory numbers look like.

\subsection{A single contract, start to finish}
Pick any transaction from the master ledger, then pull its full cash-flow
schedule:

\begin{center}
\ttfamily\small
SELECT TOP 20 * FROM dbo.transactions WHERE product\_code = '1000';\\
SELECT * FROM cf.products WHERE schedule\_id = (\\
\ \ SELECT schedule\_id FROM schemat.loans WHERE transaction\_id = 'REPLACE\_ME'\\
) ORDER BY cf\_start\_dt;
\end{center}

\begin{sqlplaceholder}[7.5cm]{Screenshot 1 --- contract $\rightarrow$ cash-flow schedule}
\normalfont\itshape\color{labgrey}
[ paste your SQL Server Management Studio result grid here ---\\
 one contract row, then its full \texttt{cf.products} schedule ]
\end{sqlplaceholder}

\subsection{Final regulatory-style output}
\begin{center}
\ttfamily\small
SELECT * FROM results.irrbb\_report ORDER BY currency, scenario;\\
SELECT * FROM results.lcr\_nsfr ORDER BY currency;
\end{center}

\begin{sqlplaceholder}[7.5cm]{Screenshot 2 --- Supervisory Outlier Test \& LCR/NSFR}
\normalfont\itshape\color{labgrey}
[ paste your SQL Server Management Studio result grid here ---\\
 \texttt{results.irrbb\_report} and/or \texttt{results.lcr\_nsfr} ]
\end{sqlplaceholder}

\begin{pitchbox}[Try it on your own bank]
Point \texttt{balance\_generate/input\_data/bank\_data.xlsx} at your own product
mix, run the \texttt{labbank\_data\_job} in Dagster, and every table above
repopulates with your numbers --- same schema, same queries, your balance sheet.
\end{pitchbox}

% =============================================================================
\section{Balance Sheet \& Product Universe}
% =============================================================================

The demo book below is 10.0 B PLN on each side, product-by-product:

\begin{center}
\small
\begin{minipage}[t]{0.47\textwidth}
\centering
\textbf{Assets} \quad(10.00 B PLN)\\[4pt]
\begin{tabular}{@{}lrr@{}}
\toprule
Product & B PLN & \% \\
\midrule
Mortgage (floating)          & 2.80  & 28.0\% \\
Consumer loan (floating)     & 1.80  & 18.0\% \\
Bond / securities (floating) & 1.60  & 16.0\% \\
Bond / securities (fixed)    & 1.20  & 12.0\% \\
Mortgage (fixed)             & 1.00  & 10.0\% \\
SME investment loan (floating) & 0.60 & 6.0\% \\
Consumer loan (fixed)        & 0.60  & 6.0\% \\
Cash / central bank reserves & 0.20  & 2.0\% \\
Interbank placement          & 0.20  & 2.0\% \\
\midrule
\textbf{Total assets} & \textbf{10.00} & \textbf{100.0\%} \\
\bottomrule
\end{tabular}
\end{minipage}\hfill
\begin{minipage}[t]{0.47\textwidth}
\centering
\textbf{Liabilities + Equity} \quad(10.00 B PLN)\\[4pt]
\begin{tabular}{@{}lrr@{}}
\toprule
Product & B PLN & \% \\
\midrule
Current account & 3.80 & 38.0\% \\
Term deposit     & 2.20 & 22.0\% \\
Savings account  & 1.60 & 16.0\% \\
Issued bond      & 1.20 & 12.0\% \\
Tier~1 equity    & 1.20 & 12.0\% \\
\midrule
\textbf{Total L + E} & \textbf{10.00} & \textbf{100.0\%} \\
\bottomrule
\end{tabular}
\end{minipage}
\end{center}

Every product is a row of parameters in an Excel sheet --- rate type, tenor,
amortisation, RWA weight, PD/LGD, LCR outflow rate, ASF/RSF weight. Adding a
new product family (credit cards, T-bills, a corporate working-capital split)
means adding a row, not writing code.

\textbf{Rate types:} \textbf{F}ixed (locked) \;|\; \textbf{V}ariable (index +
spread) \;|\; \textbf{A}dministrative (bank-set, floored at 0\%).

\subsection{How the balance sheet is generated}
Every run produces a fresh synthetic bank, not a fixed dataset --- two levers
control that: how much each product holds, and what rate each position carries.

\begin{itemize}[itemsep=2pt]
  \item \textbf{Outstanding balances} --- for each (product, client-type) bucket,
        individual transaction balances are drawn from a \textbf{truncated normal
        distribution}, parameterised per bucket (mean, std.\ dev., lower/upper
        bound) in the input Excel. Draws accumulate until their sum reaches that
        bucket's target share of the balance sheet; the final draw is trimmed so
        the bucket total matches the target exactly. Loans are drawn at their
        \textit{initial} notional and then amortised down (annuity, constant, or
        bullet, per product) to the \textit{current} outstanding balance, re-solving
        until the current-balance totals also hit target.
  \item \textbf{Interest rates} --- not one flat market rate applied everywhere:
        each position's rate is anchored to the \textbf{historical index fixing at
        its own vintage}. A fixed-rate contract looks up the fixing on its own
        \texttt{start\_date} and locks it for life, so an old fixed-rate mortgage
        still carries whatever WIBOR was years ago, exactly like a real seasoned
        book. A floating or administrative-rate product instead reprices off a
        recent fixing (as of the report date, net of its repricing delay). On top
        of the index, \texttt{client\_rt = a · index\_rt + b}, with a
        per-contract margin drawn from its own truncated-normal distribution and,
        where set, a product-level floor.
  \item \textbf{Start / maturity dates} --- start dates are random business days,
        weighted toward more recent origination (newer vintages are more likely
        than older ones, mimicking a growing back book); maturity follows from the
        product's tenor via a standard business-day calendar.
\end{itemize}

\begin{pitchbox}[Why this matters]
Because balances and rates are drawn from distributions rather than hand-set,
re-running the generator gives a structurally similar but numerically different
bank each time --- useful for testing that a metric behaves correctly across many
balance sheets, not just the one you happened to build first.
\end{pitchbox}

% =============================================================================
\section{Cash Flow Engine \& Behavioural Modelling}
% =============================================================================

For every product, the engine generates two things per period: \textbf{capital
cash flows} (contractual amortisation plus behavioural adjustments) and
\textbf{interest cash flows} (coupon payments at the contracted client rate).

\begin{itemize}[itemsep=2pt]
  \item \textbf{Loans} --- a constant prepayment rate (CPR), set per
        (product, tenor bucket) and reloaded every pipeline run --- change the
        Excel input, no code change needed.
  \item \textbf{Non-maturity deposits} --- exponential decay: a fraction leaves
        each month, the remainder is treated as sticky long-term funding.
  \item \textbf{Rate floors/caps} --- per-product, data-driven; the shipped
        demo floors every deposit-type product at 0\% so a downward shock can
        never push what the bank pays depositors negative.
\end{itemize}

The output, \texttt{cf.products}, has one row per (transaction, payment date):
contractual \textit{and} behavioural outstanding balance and cash flow, the
market forward rate, and the all-in client rate. NII is computed from interest
cash flows within the 1-year horizon; EVE uses the full behavioural run-off
schedule to maturity.

\subsection{Repricing gap \& interest rate swaps}
The \textbf{repricing gap} is the net volume of assets minus liabilities
repricing in each tenor bucket. Swaps overlay this gap \textit{without}
changing the underlying balance sheet: the swap gap is computed separately and
merged into the balance-sheet repricing gap before IRRBB is calculated --- so
LabBank can show ``balance sheet only'' vs. ``balance sheet + IRS'' side by side.

\imgfull{repricing_gap.png}

The net gap per tenor bucket --- positive means more assets than liabilities
repricing there, i.e.\ asset-sensitive. The LabBank app's Gap Analysis tab
renders this same view interactively, alongside its running cumulative total.

% =============================================================================
\section{IRRBB --- Interest Rate Risk in the Banking Book}
\label{sec:irrbb}
% =============================================================================

\imgfull[0.8\textwidth]{shock_curves.png}

Six EBA/RTS/2022/10 supervisory scenarios (parallel up/down, steepener,
flattener, short-rate up/down) plus an ``own'' scenario, applied per currency
curve, floored at 0\% per tenor.

\subsection{NII --- two components per scenario}
\begin{itemize}[itemsep=2pt]
  \item \textbf{Locked interest} --- income from cash flows already
        contracted. Immune to the shock for fixed-rate products; re-prices at
        the shocked curve (stable-margin transform) for floating-rate ones.
  \item \textbf{Renewal} --- income from capital maturing within the 1-year
        horizon, reinvested at the \textit{shocked market rate} --- this
        applies even to fixed-rate products, since a maturing fixed deposit or
        loan renews at whatever the new environment offers.
\end{itemize}

\imgfull{nii_scenarios.png}

Each bar is the full-year NII under that scenario's shocked curve; the solid
segment is locked interest, the hatched segment is renewal. Reading the whole
ladder at once shows what a single base-scenario number cannot: which shocks
help NII (parallel-up, steepener) and which hurt it (parallel-down, flattener),
and by how much relative to base.

\subsection{EVE --- present value under run-off}
\[
\mathrm{EVE} = \sum_i \mathrm{CF}_i^{\mathrm{behavioural}} \cdot d(t_i)
\qquad
\Delta\mathrm{EVE}_s = \mathrm{EVE}_s - \mathrm{EVE}_{\mathrm{base}}
\]
EVE loads every maturity (no 1-year cut-off) under a run-off assumption ---
no new business, existing book winds down. A negative $\Delta$EVE means economic
net worth falls under that scenario, typically because long fixed-rate assets
lose value faster than liabilities as rates rise --- the classic NII-vs-EVE
trade-off.

\imgfull{eve_results.png}

\subsection{Supervisory Outlier Test}
A scenario's $\Delta$EVE / $\Delta$NII gets a \textbf{50\% haircut if it comes
out positive}, before it counts toward the regulatory total:
\[
\Delta\mathrm{EVE}^{\mathrm{reg}} = \min(\Delta\mathrm{EVE}, 0)
  + 0.5 \cdot \max(\Delta\mathrm{EVE}, 0)
\]
Thresholds: EVE SOT breaches at $\Delta\mathrm{EVE}^{\mathrm{reg}}/T_1 < -15\%$;
NII SOT breaches at $\Delta\mathrm{NII}^{\mathrm{reg}}/T_1 < -5\%$.

\imgfull{eba_sot.png}

\begin{center}
\small
\begin{tabular}{@{}lrrrrc@{}}
\toprule
Scenario & $\Delta$EVE$_{\mathrm{reg}}$ (M PLN) & $/T_1$ & $\Delta$NII$_{\mathrm{reg}}$ (M PLN) & $/T_1$ & Breach \\
\midrule
Parallel +250 bps  & $-$85.4 & $-$7.11\% & $+$21.3 & $+$1.78\% & --- \\
Parallel $-$250 bps & $+$46.5 & $+$3.87\% & $-$42.5 & $-$3.54\% & --- \\
Steepener           & $-$15.9 & $-$1.32\% & $+$16.5 & $+$1.37\% & --- \\
Flattener           & $+$3.8  & $+$0.32\% & $-$29.9 & $-$2.49\% & --- \\
Short rate +350 bps & $-$29.2 & $-$2.43\% & $-$9.6  & $-$0.80\% & --- \\
Short rate $-$350 bps & $+$15.3 & $+$1.27\% & $+$4.8 & $+$0.40\% & --- \\
Own                 & $+$18.3 & $+$1.53\% & $-$17.0 & $-$1.42\% & --- \\
\bottomrule
\end{tabular}
\end{center}
\small\color{labgrey}\textit{Tier 1 capital: 1.20 B PLN. No scenario breaches either threshold for this
demo book --- the whole point of running the balance-sheet optimiser
(Section~\ref{sec:tryityourself}, roadmap) is to keep it that way as the mix changes.}
\normalsize\color{black}

% =============================================================================
\section{Liquidity Risk --- LCR \& NSFR}
% =============================================================================

\[
\mathrm{LCR} = \frac{\mathrm{HQLA}}{\text{Net outflows}_{30\mathrm{d}}} \geq 100\%
\qquad\qquad
\mathrm{NSFR} = \frac{\mathrm{ASF}}{\mathrm{RSF}} \geq 100\%
\]

\textbf{HQLA} = Level~1 liquid assets at their respective haircuts.\quad
\textbf{Net outflows} = stressed deposit run-off and interbank maturities within
30 days.\quad \textbf{ASF} = stable funding weighted by tenor and type.\quad
\textbf{RSF} = assets weighted by illiquidity.

\imgfull[0.85\textwidth]{lcr_nsfr.png}

Beyond the point-in-time ratios, the \textbf{behavioural liquidity gap} shows the
net cash position across the full tenor ladder (up to 60 months), using
behavioural cash flows --- loan prepayments included, NMD decay applied.

\imgfull{liq_gap.png}

The LabBank app's Gap Analysis tab renders this same periodic view
interactively, alongside its running cumulative total --- a sustained negative
cumulative gap is what identifies the point at which the bank would need
additional funding.

% =============================================================================
\section{Reporting \& Outputs}
% =============================================================================

LabBank produces two kinds of output, and the distinction matters: two
\textbf{static} notebook reports that read the database once and freeze the
result --- the point-in-time artefact you'd actually attach to a committee pack
or send to a regulator --- and the \textbf{interactive} Streamlit app
(Section~\ref{sec:labbankapp}) where every number recomputes as you edit. The
static reports are what you'd generate \textit{after} you've used the app to
land on a balance sheet you want to document; the app is where you get there.

\subsection{The static reports}

\textbf{ALM/Compliance report} --- \texttt{bank\_report.ipynb}, executed
against live SQL and exported to HTML/PDF via \texttt{export\_report.py}. Built
for a risk/ALM committee audience: balance sheet composition and product
breakdown, market curves (zero rates, discount factors, EBA shocked curves),
the repricing gap with and without the IRS overlay, the behavioural liquidity
gap, loan/deposit run-off profiles, CF-based NII by scenario, the EBA
Supervisory Outlier Test with the regulatory haircut applied, a linear-vs-exact
non-linearity check, LCR/NSFR, and a closing ALM-insights page (sensitivity
heatmap, repricing ladder, sensitivity tornado).

\imgfull[0.92\textwidth]{bank_report_bs_summary.png}
\begin{center}
\small\color{labgrey}\textit{Balance sheet summary from \texttt{bank\_report.ipynb} --- composition, fixed vs.\ variable rate exposure, and funding structure, all read live from \texttt{dbo.transactions}.}
\end{center}

\imgfull[0.92\textwidth]{bank_report_bs_breakdown.png}
\begin{center}
\small\color{labgrey}\textit{Balance sheet product breakdown from the same report --- every product line is a row you can trace back to \texttt{dbo.transactions}.}
\end{center}

\textbf{Finance/CFO report} --- \texttt{finance\_report.ipynb}, generated and
executed by \texttt{make\_finance\_report.py}. Same underlying cash-flow data,
reaggregated for a finance/treasury audience instead of a risk one: an
executive P\&L summary, the interest income/expense waterfall, product margin
vs.\ benchmark, funding cost structure, fixed-vs-variable exposure, the NII
bridge (locked interest vs.\ renewal), a margin profile by tenor, and Economic
Profit (margin over FTP, plus fee, minus expected loss, cost of capital, and
OpEx) both in total and per product.

\imgfull[0.92\textwidth]{finance_report_nii_waterfall.png}
\begin{center}
\small\color{labgrey}\textit{Interest income/expense waterfall and NII-by-product from \texttt{finance\_report.ipynb} --- the same NII total as the compliance report, split the way a treasury desk would ask for it.}
\end{center}

\subsection{Other outputs}

\begin{itemize}[itemsep=2pt]
  \item \textbf{Interactive exploration} --- the LabBank Streamlit app
        (Section~\ref{sec:labbankapp}) lets you stress the balance sheet, IRS
        book, and NMD assumptions and watch NII/EVE/SOT/LCR/NSFR update live,
        no SQL Server required.
  \item \textbf{Power BI \& Excel} --- \texttt{ir\_gap.pbix} connects live to
        \texttt{schemat.*}; \texttt{irrbb\_report.xlsx} and
        \texttt{lcr\_nsfr\_report.xlsx} are ready for regulatory templates.
\end{itemize}

\begin{pitchbox}[Full documentation]
This document is the pitch, not the manual. The repository ships two written
guides that go deeper than a 12-page PDF can: a \textbf{setup guide} covering
SQL Server, Python environment, and Dagster from a clean machine, and a
\textbf{methodology} document walking through every formula in this report ---
NII, EVE, the SOT haircut, LCR/NSFR, behavioural modelling --- with the exact
SQL table or Python function that computes it. Both are in \texttt{docs/} in
the repository, as Markdown and as standalone HTML.
\end{pitchbox}

% =============================================================================
\section{Why LabBank}
% =============================================================================

\begin{itemize}[itemsep=3pt]
  \item \textbf{Auditable by construction} --- one warehouse, nine schemas, one
        trail from raw transaction to regulatory output. No hidden in-memory
        state anywhere in the pipeline.
  \item \textbf{Methodology you can check} --- EBA/RTS/2022/10 shock
        construction, SOT gain haircut, behavioural CPR and NMD decay
        --- every formula has a direct SQL or Python equivalent.
  \item \textbf{Reproducible} --- Dagster enforces dependency order; any stage
        can be re-run in isolation and the run history is queryable.
  \item \textbf{Built to be extended} --- new products are Excel rows, not code
        changes; the \texttt{opt\_prep} tensor layer is already in place for the
        balance-sheet optimisation phase.
\end{itemize}

\begin{pitchbox}[Scope, honestly]
This is a laboratory project, not production banking software --- the balance
sheet is synthetic and behavioural models are illustrative rather than
calibrated to a real portfolio. What it demonstrates is \textit{how} a
production-grade ALM pipeline should be architected, end to end, and how its
metrics relate to each other without a reconciliation gap.
\end{pitchbox}

% =============================================================================
\section{Try It Yourself}
\label{sec:tryityourself}
% =============================================================================

\textbf{Fastest path --- no SQL Server, no setup beyond Python:} same four
commands as Section~\ref{sec:labbankapp} --- \texttt{streamlit run
sandbox\textbackslash app.py} after the usual venv + \texttt{pip install}.

\textbf{Full path --- your own balance sheet, with SQL Server as the audit layer:}
edit \texttt{bank\_data.xlsx}, run the \texttt{labbank\_data\_job} in Dagster,
click ``Reload my data'' in the app --- every table in
Section~\ref{sec:sqltables} repopulates with your numbers. See
Section~\ref{sec:labbankapp} for a full walkthrough of the app itself.

\subsection{Roadmap}
Three phases, one project --- the first two are shipped and are what the rest
of this document describes; the third is where LabBank is heading next.

\begin{pitchbox}[Phase 1 --- Build (shipped)]
Generate a synthetic bank from your own product mix: \texttt{balance\_generate}
$\rightarrow$ cash flows $\rightarrow$ IRRBB/liquidity, orchestrated by
Dagster (Section~\ref{sec:architecture}).
\end{pitchbox}

\begin{center}\Large$\downarrow$\end{center}

\begin{pitchbox}[Phase 2 --- Explore (shipped)]
Stress it interactively: the LabBank Streamlit sandbox
(Section~\ref{sec:labbankapp}) --- edit the balance sheet, IRS book, and NMD
assumptions and watch NII/EVE/LCR/NSFR update live.
\end{pitchbox}

\begin{center}\Large$\downarrow$\end{center}

\begin{pitchbox}[Phase 3 --- Optimise (in progress)]
Find the best balance sheet, not just stress a given one ---
\texttt{bs\_optimization/}: a deterministic constrained solver (economic
profit vs.\ $\Delta$EVE/$\Delta$NII breach trade-off), a joint
balance-sheet + swap-overlay solver, and a stochastic Monte-Carlo solver are
functional today, not yet part of the guided path. An AI-driven optimiser is
the further direction from here --- not yet started.
\end{pitchbox}

% =============================================================================
\section*{Contact}
\addcontentsline{toc}{section}{Contact}
% =============================================================================

\begin{center}
\large\color{labdark}\textbf{LabBank}\\[4pt]
\normalsize\color{labgrey}Laboratory Banking ALM System\\[16pt]
\normalsize\color{black}
\href{https://github.com/dzimius/LabBank}{github.com/dzimius/LabBank}\\[6pt]
\href{mailto:REPLACE_WITH_YOUR_EMAIL}{REPLACE\_WITH\_YOUR\_EMAIL}\\[6pt]
\href{https://www.linkedin.com/in/REPLACE_WITH_YOUR_HANDLE}{linkedin.com/in/REPLACE\_WITH\_YOUR\_HANDLE}
\end{center}

\end{document}
